% ============================================================
%  TFM — Estudio Comparativo del Espacio Latente de Modelos
%         Generativos de Texto-Imagen
%  Autor: Víctor Menéndez Ríos
%  Fragmento: Capítulo 4 — Resultados  (compilable de forma independiente)
% ============================================================
\documentclass[12pt, a4paper, twoside]{book}

% --- Encoding & Language -----------------------------------
\usepackage[utf8]{inputenc}
\usepackage[T1]{fontenc}
\usepackage[spanish]{babel}

% --- Typography --------------------------------------------
\usepackage{lmodern}
\usepackage{microtype}
\usepackage{setspace}
\onehalfspacing

% --- Maths -------------------------------------------------
\usepackage{amsmath, amssymb, amsthm}
\usepackage{bm}

% --- Layout ------------------------------------------------
\usepackage[top=3cm, bottom=3cm, inner=3.5cm, outer=2.5cm]{geometry}
\usepackage{emptypage}
\pagestyle{plain}

% --- Figures & Tables --------------------------------------
\usepackage{graphicx}
\usepackage{booktabs}
\usepackage{multirow}
\usepackage{array}
\usepackage{longtable}
\usepackage{caption}
\usepackage{subcaption}
\usepackage{float}
\usepackage{xcolor}
\usepackage{siunitx}

% --- Links -------------------------------------------------
\usepackage[hidelinks]{hyperref}
\usepackage{url}

% --- Listas ------------------------------------------------
\usepackage{enumitem}

% --- Bibliografía (natbib autor-año) -----------------------
\usepackage[authoryear, round, sort]{natbib}

% --- Custom commands ---------------------------------------
\newcommand{\z}{\mathbf{z}}
\newcommand{\x}{\mathbf{x}}
\newcommand{\eps}{\boldsymbol{\epsilon}}
\newcommand{\R}{\mathbb{R}}
\newcommand{\E}{\mathbb{E}}
\newcommand{\norm}[1]{\left\|#1\right\|_2}
\newcommand{\TODO}[1]{\textcolor{red}{\textbf{[PENDIENTE: #1]}}}

% ===========================================================
\begin{document}
\mainmatter
% ===========================================================

% ===========================================================
%  CAPÍTULO 4 — RESULTADOS
% ===========================================================
\chapter{Resultados}
\label{ch:resultados}

El capítulo organiza los resultados en cuatro bloques que siguen el orden experimental del trabajo.
La sección~\ref{sec:res:modelos} reporta la verificación de los tres modelos y las estadísticas
descriptivas del espacio latente $\z_0$. La sección~\ref{sec:res:corpus} documenta el corpus
generado. Las secciones~\ref{sec:res:geometrico} y \ref{sec:res:semantico} presentan el análisis
geométrico y semántico del pretrain. La sección~\ref{sec:res:finetune} cierra el capítulo con los
resultados del fine-tuning sobre WikiArt y el análisis de deriva latente.

% -----------------------------------------------------------
\section{Verificación de modelos y análisis del espacio latente $\z_0$}
\label{sec:res:modelos}
% -----------------------------------------------------------

\subsection{Carga y estadísticas descriptivas}

Los tres modelos se cargaron correctamente en la GPU T4 en precisión \texttt{fp16}.
La excepción es el VAE de SDXL, que requiere \texttt{float32} para evitar desbordamiento
numérico durante la decodificación, motivado por los valores de activación de mayor magnitud
del codificador de mayor resolución~\citep{podell2023sdxl}.
La Tabla~\ref{tab:estadisticas_z0} recoge las estadísticas descriptivas del latente $\z_0$
para el mismo prompt y semilla en los tres modelos.

\begin{table}[htb]
\centering
\caption{Estadísticas descriptivas del espacio latente $\z_0$ (mismo prompt y semilla).
$\bar{z}$: media; $\sigma_z$: desviación típica; $\norm{\z}$: norma L2; MAE~VAE: error
absoluto medio del ciclo codificación\textrightarrow{}decodificación en escala de píxel $[0,255]$.}
\label{tab:estadisticas_z0}
\small
\begin{tabular}{lrrrrr}
\toprule
\textbf{Modelo} & \textbf{Dim.} & $\bar{z}$ & $\sigma_z$ & $\norm{\z}$ & \textbf{MAE VAE} \\
\midrule
SD 1.5 & 16.384 & $-0{,}083$ & $0{,}856$ & $110{,}0$ & $4{,}00$ \\
SD 2.1 & 36.864 & $-0{,}162$ & $0{,}669$ & $132{,}2$ & $2{,}64$ \\
SDXL   & 65.536 & $-0{,}188$ & $0{,}595$ & $159{,}7$ & $1{,}80$ \\
\bottomrule
\end{tabular}
\end{table}

La norma L2 crece monótonamente a pesar de que la desviación típica decrece, coherente con
la expresión $\E[\norm{\z}] \approx \sqrt{N(\mu^2 + \sigma^2)}$ bajo distribución aproximadamente
gaussiana, confirmada visualmente por los histogramas de cada canal.
La mejora progresiva del MAE del ciclo encode\textrightarrow{}decode ($4{,}00 \to 2{,}64 \to 1{,}80$)
indica que los VAE de los modelos más recientes poseen mayor capacidad; el error del VAE de SDXL
(0,71\% relativo respecto al rango $[0,255]$) es subperceptual y valida la idoneidad del espacio
$\z_0$ como representación de análisis.

\subsection{Estructura espacial de los canales latentes}

\begin{figure}[htb]
\centering
\includegraphics[width=\textwidth]{data/fase1/fase1_latent_channels}
\caption{Visualización de los cuatro canales del tensor $\z_0$ para SD~1.5 (izquierda), SD~2.1
(centro) y SDXL (derecha) con el mismo prompt y semilla. El rango de valores aumenta de
$[-6,+4]$ en SD~1.5 a $[-10,+7{,}5]$ en SD~2.1.}
\label{fig:canales}
\end{figure}

La inspección visual (Figura~\ref{fig:canales}) revela diferencias cualitativas entre modelos.
En SD~1.5, los cuatro canales muestran la silueta del objeto principal con complementariedad
sistemática entre pares. En SD~2.1, la estructura semántica se conserva con mayor rango de valores,
reflejo del condicionamiento más expresivo del codificador de mayor dimensión. En SDXL, destaca
la presencia de bandas verticales periódicas que corresponden al patrón estructural del fondo de la
imagen generada: el VAE de SDXL codifica los patrones espaciales del fondo con una fidelidad
comparable al objeto principal, evidenciando una representación más distribuida y coherente con su
menor MAE.

\subsection{Verificación del ciclo de reconstrucción}

\begin{figure}[htb]
\centering
\begin{subfigure}[b]{0.32\textwidth}
    \includegraphics[width=\textwidth]{data/fase1/sd15/roundtrip}
    \caption*{SD~1.5 --- MAE $= 4{,}00$}
\end{subfigure}
\hfill
\begin{subfigure}[b]{0.32\textwidth}
    \includegraphics[width=\textwidth]{data/fase1/sd21/roundtrip}
    \caption*{SD~2.1 --- MAE $= 2{,}64$}
\end{subfigure}
\hfill
\begin{subfigure}[b]{0.32\textwidth}
    \includegraphics[width=\textwidth]{data/fase1/sdxl/roundtrip}
    \caption*{SDXL --- MAE $= 1{,}80$}
\end{subfigure}
\caption{Comparativa del ciclo codificación\textrightarrow{}decodificación (roundtrip) para los
tres modelos. La imagen generada se codifica a $\z_0$ y se decodifica de nuevo; la diferencia
absoluta media confirma la idoneidad del VAE como representación fiel del contenido visual.}
\label{fig:roundtrip}
\end{figure}

La Figura~\ref{fig:roundtrip} muestra el ciclo de reconstrucción completo para cada modelo.
Las diferencias visibles son mínimas en SDXL y apreciables solo en texturas de alta frecuencia
en SD~1.5. La reducción del error en un factor 2,2 entre SD~1.5 y SDXL confirma que la
mayor dimensión latente de SDXL ($4\times128\times128$ frente a $4\times64\times64$) preserva
más información de alta frecuencia, con independencia del proceso de difusión.

% -----------------------------------------------------------
\section{Corpus de vectores latentes}
\label{sec:res:corpus}
% -----------------------------------------------------------

Se generaron 450 vectores latentes en total (150 por modelo), almacenados como tensores
\texttt{float32} en formato $[N, 1, C, H, W]$ con los mismos prompts, semillas e hiperparámetros
aplicados a los tres modelos de forma estrictamente sincronizada.
La verificación de reproducibilidad confirmó que la regeneración de un subconjunto de muestras
con los mismos hiperparámetros produce vectores idénticos bit a bit, validando la determinismo
del muestreo DDIM~\citep{song2021ddim} con semillas fijas.

% -----------------------------------------------------------
\section{Análisis geométrico del espacio latente}
\label{sec:res:geometrico}
% -----------------------------------------------------------

\subsection{Separabilidad intra/inter-prompt}
\label{subsec:res:separabilidad}

\begin{table}[htb]
\centering
\caption{Estadísticas de distancias L2 intra e inter-prompt (60 latentes base por modelo).
$\mu_\text{intra}$: distancia media entre latentes del mismo prompt con semillas distintas;
$\mu_\text{inter}$: distancia media entre latentes de prompts distintos;
sep $= \mu_\text{inter} / \mu_\text{intra}$; KS: estadístico de Kolmogorov--Smirnov bilateral;
$p$: nivel de significación.}
\label{tab:separabilidad}
\small
\begin{tabular}{lrrrrrrr}
\toprule
\textbf{Modelo} & $\mu_\text{intra}$ & $\sigma_\text{intra}$ & $\mu_\text{inter}$ &
$\sigma_\text{inter}$ & \textbf{sep} & \textbf{KS} & \textbf{$p$} \\
\midrule
SD 1.5 & 151,8 & 16,4 & 148,1 & 22,4 & 0,976 & 0,142 & 0,173 \\
SD 2.1 & 195,8 & 22,4 & 189,1 & 30,7 & 0,966 & 0,187 & 0,031 \\
SDXL   & 230,5 & 35,0 & 257,6 & 42,9 & \textbf{1,118} & 0,304 & $<{0{,}001}$ \\
\bottomrule
\end{tabular}
\end{table}

\begin{figure}[htbp]
\centering
\begin{subfigure}[b]{\textwidth}
    \includegraphics[width=\textwidth]{data/fase3/results/geo_sd15_distancias}
    \caption*{SD~1.5}
\end{subfigure}\\[0.4em]
\begin{subfigure}[b]{\textwidth}
    \includegraphics[width=\textwidth]{data/fase3/results/geo_sd21_distancias}
    \caption*{SD~2.1}
\end{subfigure}\\[0.4em]
\begin{subfigure}[b]{\textwidth}
    \includegraphics[width=\textwidth]{data/fase3/results/geo_sdxl_distancias}
    \caption*{SDXL}
\end{subfigure}
\caption{Distribuciones de distancias L2 intra-prompt e inter-prompt. En SD~1.5 y SD~2.1 los
histogramas se superponen casi completamente (sep $< 1$); en SDXL están claramente separados
(sep $> 1$, $p < 0{,}001$).}
\label{fig:distancias}
\end{figure}

SD~1.5 y SD~2.1 presentan sep $< 1$ (0,976 y 0,966), indicando que los latentes del mismo
prompt con distintas semillas están en promedio \emph{más alejados} que latentes de prompts
distintos (Tabla~\ref{tab:separabilidad}, Figura~\ref{fig:distancias}). El test KS no rechaza
la igualdad de distribuciones para SD~1.5 ($p = 0{,}173$) y solo marginalmente para SD~2.1
($p = 0{,}031$). En estos modelos, la variabilidad estocástica del proceso de denoising
desplaza el latente tanto como un cambio de contenido semántico completo: la relación
señal-ruido es desfavorable para las técnicas de edición basadas en aritmética de vectores
en $\z_0$.

SDXL, en cambio, presenta sep $= 1{,}118$ con $p < 0{,}001$: los latentes del mismo prompt
están sistemáticamente más próximos que los de prompts distintos. El contenido del prompt
es el predictor dominante de la posición del latente, con la variabilidad estocástica
subordinada al significado semántico.

\subsection{Dimensionalidad efectiva e inspección PCA}
\label{subsec:res:dimensionalidad}

Los valores propios de la covarianza se obtienen mediante SVD thin de la matriz de datos
centrada sobre los 60 latentes base aplanados a vectores de dimensión $D$ (16.384, 36.864
o 65.536 elementos según el modelo). La Tabla~\ref{tab:dimensionalidad} reúne el participation
ratio (PR), los umbrales de varianza acumulada y las fracciones explicadas por las dos primeras
componentes, evitando así la repetición de información entre tablas.

\begin{table}[htb]
\centering
\caption{Dimensionalidad efectiva e inspección de los primeros componentes PCA (60 latentes
base por modelo). PR: participation ratio; dims@80/90/95\%: número de componentes para
alcanzar el umbral de varianza acumulada indicado; PC1, PC2: fracción de varianza explicada
por el primero y segundo componente respectivamente.}
\label{tab:dimensionalidad}
\small
\begin{tabular}{lrrrrrrr}
\toprule
\textbf{Modelo} & \textbf{Dim.\ nom.} & \textbf{PR} & \textbf{dims@80\%} &
\textbf{dims@90\%} & \textbf{dims@95\%} & \textbf{PC1 (\%)} & \textbf{PC2 (\%)} \\
\midrule
SD 1.5 & 16.384 & 12,35 & 25 & 38 & 46 & 19,2 & 17,7 \\
SD 2.1 & 36.864 & 12,33 & 27 & 39 & 47 & 20,1 & 17,1 \\
SDXL   & 65.536 & 11,56 & 22 & 35 & 45 & 19,1 & 14,6 \\
\bottomrule
\end{tabular}
\end{table}

\begin{figure}[htb]
\centering
\includegraphics[width=0.85\textwidth]{data/fase3/results/snapshot_cross_model_pca}
\caption{Varianza acumulada en función del número de componentes principales para los tres
modelos. Las curvas son prácticamente indistinguibles hasta el componente 10, donde SDXL
diverge ligeramente hacia una mayor concentración.}
\label{fig:pca_varianza}
\end{figure}

El participation ratio PR $\approx 12$ es prácticamente idéntico en los tres modelos pese
a que sus dimensiones nominales difieren en un factor de cuatro
(Tabla~\ref{tab:dimensionalidad}, Figura~\ref{fig:pca_varianza}).
El número de componentes para capturar el 80\% de la varianza (22--27) representa
menos del 0,04\% de la dimensión nominal de SDXL, confirmando que la baja dimensionalidad
intrínseca es una propiedad arquitectónica compartida, no vinculada al tamaño del espacio.

La diferencia más notable reside en la tercera componente de SDXL, que explica el 13,3\%
de la varianza (acumulado del 47,1\% con tres componentes), frente al 6,1\% y 5,3\% en los
modelos anteriores. Esa tercera componente coincide con la dirección semántica de
iluminación (véase la sección~\ref{sec:res:semantico}).

\subsection{Correlación entre distancia latente y similitud de imagen}
\label{subsec:res:correlacion}

Se muestrearon aleatoriamente 500 pares de los $\binom{60}{2} = 1.770$ posibles (semilla fija 42)
y se calcularon las correlaciones de Pearson ($r$) y Spearman ($\rho$) entre la distancia L2
latente y las métricas de disimilitud MSE y $-\text{SSIM}$~\citep{wang2004ssim}.

\begin{table}[htb]
\centering
\caption{Correlación entre distancia latente L2 y métricas de disimilitud de imagen
(500 pares por modelo, semilla 42).}
\label{tab:correlacion}
\small
\begin{tabular}{lrrrr}
\toprule
\textbf{Modelo} & $r(\text{L2, MSE})$ & $\rho(\text{L2, MSE})$ &
$r(\text{L2}, -\text{SSIM})$ & \textbf{$p$-valor} \\
\midrule
SD 1.5 & 0,740 & 0,763 & 0,796 & $\ll 10^{-80}$ \\
SD 2.1 & 0,718 & 0,753 & 0,761 & $\ll 10^{-80}$ \\
SDXL   & 0,680 & 0,698 & 0,691 & $\ll 10^{-68}$ \\
\bottomrule
\end{tabular}
\end{table}

La correlación es positiva, fuerte y altamente significativa en los tres modelos
(Tabla~\ref{tab:correlacion}). El orden $r_{\text{sd15}} > r_{\text{sd21}} > r_{\text{sdxl}}$
sugiere que el VAE de menor resolución de SD~1.5 actúa como filtro de información de alta
frecuencia más agresivo, produciendo latentes cuya distancia L2 está más alineada con
métricas globales como el MSE. El mayor MAE del ciclo roundtrip de SD~1.5 y la mayor
correlación con MSE son dos caras del mismo fenómeno: el VAE de SD~1.5 descarta más
información de alta frecuencia, haciendo que las distancias latentes reflejen más fielmente
las diferencias semánticas globales capturadas por el MSE.

% -----------------------------------------------------------
\section{Análisis semántico del espacio latente}
\label{sec:res:semantico}
% -----------------------------------------------------------

\subsection{Discriminabilidad de las direcciones semánticas}
\label{subsec:res:discriminabilidad}

La dirección semántica de cada grupo de atributo $g$ se define como la diferencia de medias
(Ec.~\ref{eq:dir_semantica}, §3.4). La Tabla~\ref{tab:normas_dir} recoge las normas medias
por tipo de atributo y modelo.

\begin{table}[htb]
\centering
\caption{Norma media ($\pm$ desviación típica) de las direcciones semánticas
$\mathbf{d}_g = \bar{\z}(g,v^+) - \bar{\z}(g,v^-)$ por tipo de atributo y modelo.}
\label{tab:normas_dir}
\small
\begin{tabular}{lrrr}
\toprule
\textbf{Tipo de atributo} & \textbf{SD 1.5} & \textbf{SD 2.1} & \textbf{SDXL} \\
\midrule
Color (cálido/frío)            & $65{,}4 \pm 4{,}1$  & $90{,}4 \pm 10{,}9$ & $183{,}4 \pm 18{,}8$ \\
Iluminación (brillante/oscuro) & $55{,}8 \pm 11{,}4$ & $76{,}1 \pm 13{,}5$ & $149{,}8 \pm 29{,}1$ \\
Estilo (foto/artístico)        & $51{,}6 \pm 15{,}0$ & $64{,}8 \pm 12{,}3$ & $106{,}5 \pm 32{,}7$ \\
\bottomrule
\end{tabular}
\end{table}

\begin{figure}[htbp]
\centering
\begin{subfigure}[b]{\textwidth}
    \includegraphics[width=\textwidth]{data/fase4/results/sem_sd15_direcciones}
    \caption*{SD~1.5}
\end{subfigure}\\[0.4em]
\begin{subfigure}[b]{\textwidth}
    \includegraphics[width=\textwidth]{data/fase4/results/sem_sd21_direcciones}
    \caption*{SD~2.1}
\end{subfigure}\\[0.4em]
\begin{subfigure}[b]{\textwidth}
    \includegraphics[width=\textwidth]{data/fase4/results/sem_sdxl_direcciones}
    \caption*{SDXL}
\end{subfigure}
\caption{Norma media de las direcciones semánticas por tipo de atributo y modelo.
Las barras de error muestran la desviación típica entre los cinco grupos de cada tipo.}
\label{fig:normas}
\end{figure}

Se observan dos patrones consistentes (Tabla~\ref{tab:normas_dir}, Figura~\ref{fig:normas}).
La \textbf{jerarquía interna} color $>$ iluminación $>$ estilo se mantiene en los tres modelos.
Las normas de SDXL son 2,1--2,8$\times$ las de SD~1.5, un crecimiento proporcionalmente
mayor que el de las normas de los latentes base (factor $\sim 1{,}8\times$), indicando mayor
expresividad semántica neta del doble condicionamiento textual.

\subsection{Consistencia intra-tipo de las direcciones semánticas}
\label{subsec:res:consistencia}

\begin{table}[htb]
\centering
\caption{Similitud coseno media intra-tipo ($\pm$ desviación típica entre pares de grupos).
Valores en $[-1, 1]$; mayor indica mayor universalidad (escena-independencia) de la dirección.}
\label{tab:consistencia}
\small
\begin{tabular}{lrrr}
\toprule
\textbf{Tipo} & \textbf{SD 1.5} & \textbf{SD 2.1} & \textbf{SDXL} \\
\midrule
Color       & $0{,}300 \pm 0{,}061$ & $0{,}229 \pm 0{,}112$ & $\mathbf{0{,}707 \pm 0{,}048}$ \\
Iluminación & $0{,}109 \pm 0{,}091$ & $0{,}155 \pm 0{,}062$ & $\mathbf{0{,}517 \pm 0{,}111}$ \\
Estilo      & $0{,}028 \pm 0{,}081$ & $0{,}010 \pm 0{,}082$ & $\mathbf{0{,}209 \pm 0{,}081}$ \\
\bottomrule
\end{tabular}
\end{table}

\begin{figure}[htbp]
\centering
\begin{subfigure}[b]{\textwidth}
    \includegraphics[width=\textwidth]{data/fase4/results/sem_sd15_angulos}
    \caption*{SD~1.5}
\end{subfigure}\\[0.4em]
\begin{subfigure}[b]{\textwidth}
    \includegraphics[width=\textwidth]{data/fase4/results/sem_sd21_angulos}
    \caption*{SD~2.1}
\end{subfigure}\\[0.4em]
\begin{subfigure}[b]{\textwidth}
    \includegraphics[width=\textwidth]{data/fase4/results/sem_sdxl_angulos}
    \caption*{SDXL}
\end{subfigure}
\caption{Mapa de calor de similitudes coseno entre las quince direcciones semánticas
(cinco por tipo) para cada modelo. En SDXL los cinco grupos de color presentan alta
similitud mutua (bloque rojo superior izquierdo), evidenciando una dirección de color
escena-independiente.}
\label{fig:consistencia}
\end{figure}

SD~1.5 y SD~2.1 muestran consistencias bajas en todos los tipos: las direcciones semánticas
no son universales sino dependientes del contexto de la escena
(Tabla~\ref{tab:consistencia}, Figura~\ref{fig:consistencia}).
SDXL muestra consistencias notablemente más altas: 0,707 para color (ángulo medio de ${\sim}45^{\circ}$
entre pares), 0,517 para iluminación y 0,209 para estilo. En SDXL existe una \textit{dirección
universal del color} escena-independiente con intensidad equiparable a la del espacio $\mathcal{W}$
de StyleGAN para atributos faciales~\citep{shen2020interfacegan}.

Destaca la ausencia de mejora entre SD~1.5 y SD~2.1 pese al mayor tamaño del codificador
OpenCLIP ViT-H/14: la consistencia de estilo retrocede incluso de 0,028 a 0,010. El salto
cualitativo se produce únicamente con el doble codificador de SDXL.

\subsection{Ortogonalidad inter-tipo (disentanglement)}
\label{subsec:res:ortogonalidad}

\begin{table}[htb]
\centering
\caption{Similitud coseno absoluta media inter-tipo ($\pm$ desviación típica).
Valores en $[0, 1]$; menor indica mejor disentanglement entre tipos de atributo.}
\label{tab:ortogonalidad}
\small
\begin{tabular}{lrrr}
\toprule
\textbf{Par de tipos} & \textbf{SD 1.5} & \textbf{SD 2.1} & \textbf{SDXL} \\
\midrule
Color vs Iluminación  & $0{,}093 \pm 0{,}048$ & $0{,}083 \pm 0{,}052$ & $0{,}119 \pm 0{,}086$ \\
Color vs Estilo       & $0{,}056 \pm 0{,}050$ & $0{,}051 \pm 0{,}039$ & $0{,}227 \pm 0{,}121$ \\
Iluminación vs Estilo & $0{,}068 \pm 0{,}058$ & $0{,}087 \pm 0{,}079$ & $0{,}328 \pm 0{,}130$ \\
\bottomrule
\end{tabular}
\end{table}

Se registra una \textbf{tensión entre consistencia y disentanglement}: SD~1.5 y SD~2.1
presentan excelente ortogonalidad ($|\cos| < 0{,}10$) pero sus direcciones son
contexto-dependientes; SDXL invierte el equilibrio con mayor consistencia pero mayor
entrelazamiento, especialmente entre iluminación y estilo ($|\cos| = 0{,}328$)
(Tabla~\ref{tab:ortogonalidad}).

El entrelazamiento iluminación-estilo en SDXL puede reflejar correlaciones visuales genuinas
del dominio artístico: los estilos plásticos poseen tratamientos lumínicos característicos
(claroscuro barroco, luz difusa impresionista, contraste saturado expresionista), cuya
captura semántica coherente se traduciría necesariamente en correlación entre ambas
direcciones.

\subsection{Alineación de las direcciones semánticas con la base PCA}
\label{subsec:res:energia_pca}

\begin{table}[htb]
\centering
\caption{Fracción media de energía de las direcciones semánticas en los diez primeros
componentes principales ($f_g$, media sobre los cinco grupos del tipo).}
\label{tab:energia_pca}
\small
\begin{tabular}{lrrr}
\toprule
\textbf{Tipo} & \textbf{SD 1.5} & \textbf{SD 2.1} & \textbf{SDXL} \\
\midrule
Color       & 0,72 & 0,67 & \textbf{0,95} \\
Iluminación & 0,56 & 0,50 & \textbf{0,93} \\
Estilo      & 0,51 & 0,51 & \textbf{0,79} \\
\bottomrule
\end{tabular}
\end{table}

En SDXL, el 95\% de la energía de las direcciones de color y el 93\% de las de iluminación
se concentra en los diez primeros componentes principales (Tabla~\ref{tab:energia_pca}).
El PC0 está sistemáticamente alineado con las direcciones de color y el PC3 con las de
iluminación, permitiendo interpretar esas componentes como \textit{dimensiones de color} y
\textit{dimensión de iluminación} del espacio latente de SDXL. Esta propiedad establece
un puente directo con GANSpace~\citep{harkonen2020ganspace}: el PCA sobre $\z_0$ de SDXL
produce componentes semánticamente interpretables sin supervisión, con la ventaja adicional
de que $\z_0$ es directamente accesible mediante inversión DDIM~\citep{song2021ddim} sin
ninguna red de mapeo explícita.

En SD~1.5 y SD~2.1, la fracción de energía (50--72\%) es insuficiente para una alineación
estable; las componentes PCA reflejan principalmente varianza estructural de las escenas,
no organización de atributos visuales. El umbral a partir del cual el PCA descubre semántica
parece situarse entre SD~2.1 y SDXL, coincidiendo con el salto en el condicionamiento textual.

% -----------------------------------------------------------
\section{Fine-tuning y deriva del espacio latente}
\label{sec:res:finetune}
% -----------------------------------------------------------

Los tres modelos se sometieron a fine-tuning LoRA sobre el subconjunto de 500 imágenes de
WikiArt~\citep{saleh2015wikiart}, con los mismos hiperparámetros y precisión, adaptando
únicamente el número de pasos de SDXL (800 frente a 1.000) a las restricciones de
memoria de la GPU T4.
Tras el fine-tuning se regeneró el corpus completo con cada modelo adaptado,
comparando los latentes resultantes ($\z_0^{\text{ft}}$) con los del pretrain ($\z_0^{\text{pre}}$)
para los mismos prompts y semillas. La congelación del VAE garantiza que ambos conjuntos sean
comparables en el mismo espacio vectorial.

\subsection{Deriva geométrica}
\label{subsec:res:deriva_geometrica}

\begin{table}[htb]
\centering
\caption{Métricas de deriva geométrica pre/post fine-tuning WikiArt. $\bar{\delta} =
\|\z_0^\text{ft} - \z_0^\text{pre}\|_2$ medio por latente; $\cos_c$: similitud coseno entre
centroides de la nube pre y post; KS: estadístico del test Kolmogorov--Smirnov sobre las
distribuciones de distancias inter-prompt ($p$ entre paréntesis).}
\label{tab:deriva_geometrica}
\small
\begin{tabular}{lrrrrrrrr}
\toprule
\textbf{Modelo} & \textbf{sep\textsubscript{pre}} & \textbf{sep\textsubscript{ft}} &
$\Delta\text{sep}$ & \textbf{PR\textsubscript{pre}} & \textbf{PR\textsubscript{ft}} &
$\Delta\text{PR}$ & $\cos_c$ & $\bar{\delta}$ \\
\midrule
SD 1.5 & 0,976 & 0,974 & $-0{,}001$ & 12,35 & 11,68 & $-0{,}67$ & 0,991 & \phantom{00}36,6 \\
SD 2.1 & 0,966 & 0,959 & $-0{,}007$ & 12,33 & 10,75 & $-1{,}58$ & 0,993 & \phantom{00}39,6 \\
SDXL   & 1,118 & 1,013 & $\mathbf{-0{,}105}$ & 11,56 & 10,11 & $-1{,}45$ & 0,914 &
$\mathbf{138,4}$ \\
\bottomrule
\end{tabular}
\end{table}

\begin{figure}[htbp]
\centering
\includegraphics[width=\textwidth]{data/fase5/deriva_sd15}
\caption{Deriva latente pre/post fine-tuning WikiArt para SD~1.5. La distribución de
distancias, la separabilidad y el PR apenas varían; $\bar{\delta} = 36{,}6$. La consistencia
semántica y la energía PCA de las direcciones se mantienen estables. La base PCA conserva
una alineación de Grassmann de 0,82 en los diez primeros componentes.}
\label{fig:deriva_sd15}
\end{figure}

\begin{figure}[htbp]
\centering
\includegraphics[width=\textwidth]{data/fase5/deriva_sd21}
\caption{Deriva latente pre/post fine-tuning WikiArt para SD~2.1. La geometría es
prácticamente estable. La consistencia de estilo muestra el único incremento apreciable
entre todos los modelos ($0{,}010 \to 0{,}023$, $+130\%$). La alineación PCA top-10
es la más alta de los tres modelos (0,888).}
\label{fig:deriva_sd21}
\end{figure}

\begin{figure}[htbp]
\centering
\includegraphics[width=\textwidth]{data/fase5/deriva_sdxl}
\caption{Deriva latente pre/post fine-tuning WikiArt para SDXL. El rango del histograma
de desplazamiento por latente alcanza $\sim 250$ frente a $\sim 100$ en SD~1.5 y SD~2.1.
La separabilidad cae de 1,118 a 1,013. La consistencia semántica disminuye en los tres
atributos (color: $-22\%$, iluminación: $-52\%$, estilo: $-37\%$). La alineación PCA
top-10 cae a 0,537.}
\label{fig:deriva_sdxl}
\end{figure}

SD~1.5 y SD~2.1 muestran deriva geométrica prácticamente nula (Tabla~\ref{tab:deriva_geometrica},
Figuras~\ref{fig:deriva_sd15}--\ref{fig:deriva_sd21}).
El ratio de separabilidad varía en menos de 0,01 puntos; los desplazamientos medios por latente
($\bar{\delta} = 36{,}6$ y $39{,}6$) representan el 24\% y el 20\% de las distancias intra-prompt
pretrain, respectivamente. Los centroides se mantienen casi idénticos ($\cos_c > 0{,}99$) y el
test KS no rechaza la igualdad de distribuciones de distancias inter-prompt ($p = 0{,}182$ y
$p = 0{,}665$).

SDXL presenta un perfil cualitativamente distinto (Figura~\ref{fig:deriva_sdxl}).
El desplazamiento medio por latente ($\bar{\delta} = 138{,}4$) equivale al 53\% de la distancia
intra-prompt pretrain ($\mu_\text{intra} \approx 261$): los latentes regenerados de SDXL se han
desplazado más de la mitad de la distancia típica entre semillas del mismo prompt.
El ratio de separabilidad cae 0,105 puntos, perdiendo casi toda la ventaja característica de
SDXL pretrain.
La similitud coseno entre centroides cae a 0,914, registrando la mayor desviación direccional.
A pesar de esta magnitud, el test KS tampoco rechaza la igualdad de distribuciones
($p = 0{,}120$), indicando que la \emph{forma} global de la distribución de distancias se
preserva aunque su posición cambie.

Respecto al participation ratio, SD~2.1 registra la mayor pérdida de dimensionalidad efectiva
($\Delta\text{PR} = -1{,}58$), ligeramente superior a la de SDXL ($-1{,}45$), mientras que
SD~1.5 pierde solo $-0{,}67$ dimensiones. El fine-tuning concentra la varianza en menos
dimensiones en los tres modelos, pero con intensidad mayor en SD~2.1 y SDXL.

\subsection{Rotación de la base PCA tras el fine-tuning}
\label{subsec:res:pca_alineacion}

\begin{table}[htb]
\centering
\caption{Alineación de Grassmann entre los subespacios PCA pretrain y post fine-tuning.
El valor 1 indica subespacios idénticos; valores menores indican rotación de la base.}
\label{tab:pca_alineacion}
\small
\begin{tabular}{lrr}
\toprule
\textbf{Modelo} & \textbf{Alineación top-10} & \textbf{Alineación top-50} \\
\midrule
SD 1.5 & 0,817 & 0,725 \\
SD 2.1 & \textbf{0,888} & \textbf{0,800} \\
SDXL   & 0,537 & 0,282 \\
\bottomrule
\end{tabular}
\end{table}

SD~1.5 y SD~2.1 mantienen alineaciones de 0,82--0,89 en los diez primeros componentes
(Tabla~\ref{tab:pca_alineacion}): el fine-tuning no reorganiza significativamente las
direcciones de mayor varianza.
En SDXL, la alineación cae a 0,537 en los diez primeros componentes y a 0,282 en los
cincuenta primeros: el subespacio de máxima varianza post fine-tuning comparte algo más de
la mitad de la orientación del subespacio pretrain en los componentes principales y menos
de un tercio en el subespacio extendido.

\subsection{Deriva semántica de las direcciones de atributo}
\label{subsec:res:deriva_semantica}

\begin{table}[htb]
\centering
\caption{Consistencia intra-tipo pre y post fine-tuning WikiArt. Los valores entre
paréntesis indican el cambio porcentual relativo.}
\label{tab:deriva_semantica_consistencia}
\small
\begin{tabular}{lcccccc}
\toprule
\textbf{Tipo} & \multicolumn{2}{c}{\textbf{SD 1.5}} & \multicolumn{2}{c}{\textbf{SD 2.1}} &
\multicolumn{2}{c}{\textbf{SDXL}} \\
\cmidrule(lr){2-3}\cmidrule(lr){4-5}\cmidrule(lr){6-7}
 & Pre & Post & Pre & Post & Pre & Post \\
\midrule
Color       & 0,300 & 0,276 {\small($-8\%$)}   & 0,229 & 0,231 {\small($+1\%$)}   &
0,707 & 0,554 {\small($-22\%$)} \\
Iluminación & 0,109 & 0,116 {\small($+6\%$)}   & 0,155 & 0,163 {\small($+5\%$)}   &
0,517 & 0,247 {\small($-52\%$)} \\
Estilo      & 0,028 & 0,027 {\small($-4\%$)}   & 0,010 & 0,023 {\small($+130\%$)} &
0,209 & 0,132 {\small($-37\%$)} \\
\bottomrule
\end{tabular}
\end{table}

SD~1.5 y SD~2.1 exhiben estabilidad semántica casi total: los cambios en la consistencia
intra-tipo son inferiores al 8\% relativo en todos los atributos, con la única excepción de la
consistencia de estilo en SD~2.1, que pasa de 0,010 a 0,023 ($+130\%$ relativo)
(Tabla~\ref{tab:deriva_semantica_consistencia}).
Este incremento, aunque grande en términos porcentuales, parte de un valor de base prácticamente
nulo; el resultado post fine-tuning (0,023) sigue siendo inferior en un orden de magnitud a la
consistencia de estilo pretrain de SDXL (0,209).

SDXL muestra una deriva semántica severa en todos los atributos: $-22\%$ en color, $-52\%$
en iluminación y $-37\%$ en estilo. Los valores post fine-tuning de SDXL siguen siendo
superiores a los pretrain de SD~1.5 y SD~2.1, confirmando que la representación semántica
de SDXL, aunque perturbada, mantiene ventaja cualitativa. La disminución más severa
corresponde a la iluminación, el atributo cuya dirección estaba parcialmente alineada con
el estilo antes del fine-tuning ($|\cos| = 0{,}328$, Tabla~\ref{tab:ortogonalidad}), sugiriendo
que la reorganización inducida por WikiArt afecta especialmente a las propiedades con
correlación visual en el dominio artístico.

\subsection{Deriva en la magnitud y la energía PCA de las direcciones semánticas}
\label{subsec:res:deriva_normas_energia}

Además de la consistencia, el fine-tuning modifica la magnitud de las direcciones y su
alineación con la base de varianza principal. La Tabla~\ref{tab:deriva_normas_energia}
resume la norma media y la energía PCA pre y post fine-tuning para cada tipo de atributo.

\begin{table}[htb]
\centering
\caption{Norma media y energía PCA ($f_g$ en los diez primeros componentes) de las
direcciones semánticas antes y después del fine-tuning WikiArt. Los cambios porcentuales
relativos se indican entre paréntesis. Solo se muestra SDXL para los atributos con cambio
apreciable; en SD~1.5 y SD~2.1 todos los deltas son inferiores al 7\% y se omiten.}
\label{tab:deriva_normas_energia}
\small
\begin{tabular}{lccccccc}
\toprule
 & \multicolumn{3}{c}{\textbf{Norma media ($\norm{\mathbf{d}_g}$)}} &
\multicolumn{3}{c}{\textbf{Energía PCA ($f_g$)}} \\
\cmidrule(lr){2-4}\cmidrule(lr){5-7}
\textbf{Tipo} & \textbf{Pre} & \textbf{Post} & \textbf{$\Delta$\%} &
\textbf{Pre} & \textbf{Post} & \textbf{$\Delta$\%} \\
\midrule
\multicolumn{7}{l}{\textit{SDXL}} \\
Color       & 183,4 & 129,4 & $-29\%$ & 0,948 & 0,925 & $-2\%$ \\
Iluminación & 149,8 & \phantom{1}98,7 & $-34\%$ & 0,932 & 0,845 & $-9\%$ \\
Estilo      & 106,5 & \phantom{1}90,9 & $-15\%$ & 0,794 & 0,789 & $-1\%$ \\
\addlinespace
\multicolumn{7}{l}{\textit{SD~2.1} (referencia)} \\
Estilo      & \phantom{1}64,8 & \phantom{1}60,2 & $-7\%$ & 0,511 & 0,553 & $+8\%$ \\
\bottomrule
\end{tabular}
\end{table}

En SDXL, las normas de las direcciones de color e iluminación caen un 29\% y un 34\%
respectivamente (Tabla~\ref{tab:deriva_normas_energia}). Esta reducción de magnitud acompaña
y amplifica la caída de consistencia: no solo las direcciones se vuelven menos universales,
sino que la diferencia vectorial entre los dos valores de cada atributo se acorta. La energía PCA
se mantiene relativamente estable (caída de 2--9\%), indicando que las direcciones post
fine-tuning de SDXL, aunque más cortas y menos consistentes, siguen proyectándose
principalmente sobre el subespacio de máxima varianza.

En SD~2.1, la norma de la dirección de estilo cae un 7\% pero su energía PCA mejora
ligeramente ($+8\%$), coherente con el pequeño incremento de consistencia (+130\% relativo)
observado en la Tabla~\ref{tab:deriva_semantica_consistencia}: el fine-tuning sobre WikiArt
da cierto orden semántico a la dirección de estilo de SD~2.1, aunque en magnitud
absolutamente insuficiente para alcanzar los valores de SDXL.

\subsection{Comparativa cross-modelo de la deriva latente}
\label{subsec:res:deriva_cross}

\begin{figure}[htbp]
\centering
\includegraphics[width=\textwidth]{data/deriva_cross_model}
\caption{Comparativa cross-modelo de la deriva latente tras fine-tuning WikiArt.
Fila superior (izq.\ a der.): $\Delta$sep, $\Delta$PR y alineación PCA top-10.
Fila inferior: $\Delta$consistencia intra-tipo, $\Delta$energía PCA de las direcciones
semánticas y desplazamiento latente medio $\bar{\delta}$.
El modelo SDXL (verde) domina todas las métricas de deriva geométrica y semántica.}
\label{fig:deriva_cross}
\end{figure}

La Figura~\ref{fig:deriva_cross} sintetiza la comparativa entre los tres modelos.
El patrón es uniforme y consistente: SD~1.5 y SD~2.1 presentan deltas cercanos a cero en
prácticamente todas las métricas, mientras que SDXL exhibe las mayores magnitudes de cambio:
\begin{itemize}[noitemsep, leftmargin=2em]
    \item Separabilidad: $\Delta\text{sep} = -0{,}105$ (frente a $-0{,}001$ y $-0{,}007$)
    \item Desplazamiento latente: $\bar{\delta} = 138{,}4$ ($\times3{,}5$--$3{,}8$ respecto a SD~1.5/SD~2.1)
    \item Alineación PCA top-10: 0,537 (frente a 0,817 y 0,888)
    \item $\Delta$consistencia color: $-0{,}153$; $\Delta$consistencia iluminación: $-0{,}270$
\end{itemize}

La notable excepción es $\Delta\text{PR}$: SD~2.1 registra la mayor pérdida de dimensionalidad
efectiva ($-1{,}58$), ligeramente superior a la de SDXL ($-1{,}45$), mientras que SD~1.5
pierde solo $-0{,}67$. Este patrón sugiere que SD~2.1, aunque geométricamente estable a nivel
de distancias, experimenta la mayor concentración dimensional de varianza tras el fine-tuning:
el espacio latente mantiene su posición pero se vuelve localmente más plano.

% -----------------------------------------------------------
\section*{Síntesis del capítulo}
% -----------------------------------------------------------

\noindent
Los resultados se articulan en torno a tres observaciones centrales.

\medskip
\noindent\textbf{PR$\approx 12$: convergencia geométrica universal.}
La dimensionalidad intrínseca es prácticamente idéntica en los tres modelos con independencia
de sus dimensiones nominales, confirmando que esta propiedad está fijada por el proceso de
difusión, no por la arquitectura.

\medskip
\noindent\textbf{Consistencia semántica: divergencia dependiente del codificador.}
El doble codificador de SDXL induce consistencias de atributo 2,3--21 veces superiores a las
de SD~1.5 y SD~2.1. El umbral de universalidad semántica en $\z_0$ no se cruza escalando un
único codificador CLIP.

\medskip
\noindent\textbf{Fine-tuning WikiArt: inversión del orden de fragilidad.}
SDXL, con la representación pretrain más organizada, sufre la mayor perturbación geométrica
y semántica; SD~1.5 y SD~2.1 muestran deriva prácticamente nula. La organización semántica
pretrain amplifica la susceptibilidad a la disrupción por adaptación de dominio.

% ===========================================================
%  REFERENCIAS (solo las citadas en este capítulo)
% ===========================================================
\bibliographystyle{apalike}

\begin{thebibliography}{99}

\bibitem[H\"{a}rkonen et al., 2020]{harkonen2020ganspace}
H\"{a}rkonen, E., Hertzmann, A., Lehtinen, J., y Paris, S. (2020).
GANSpace: Discovering interpretable GAN controls.
\textit{Advances in Neural Information Processing Systems, 33}.

\bibitem[Ho et al., 2020]{ho2020ddpm}
Ho, J., Jain, A., y Abbeel, P. (2020).
Denoising diffusion probabilistic models.
\textit{Advances in Neural Information Processing Systems, 33}, 6840--6851.

\bibitem[Hu et al., 2022]{hu2022lora}
Hu, E.~J., Shen, Y., Wallis, P., Allen-Zhu, Z., Li, Y., Wang, S., Wang, L., y Chen, W. (2022).
LoRA: Low-rank adaptation of large language models.
\textit{International Conference on Learning Representations}.

\bibitem[Kwon et al., 2023]{kwon2023diffusion}
Kwon, M., Jeong, J., y Uh, Y. (2023).
Diffusion models already have a semantic latent space.
\textit{International Conference on Learning Representations}.

\bibitem[Podell et al., 2023]{podell2023sdxl}
Podell, D., English, Z., Lacey, K., Blattmann, A., Dockhorn, T., M\"{u}ller, J.,
Penna, J., y Rombach, R. (2023).
SDXL: Improving latent diffusion models for high-resolution image synthesis.
\textit{arXiv}. \url{https://arxiv.org/abs/2307.01952}

\bibitem[Pope et al., 2021]{pope2021intrinsic}
Pope, P., Zhu, C., Abdelkader, A., Goldblum, M., y Goldstein, T. (2021).
The intrinsic dimension of images and its impact on learning.
\textit{International Conference on Learning Representations}.

\bibitem[Rombach et al., 2022]{rombach2022ldm}
Rombach, R., Blattmann, A., Lorenz, D., Esser, P., y Ommer, B. (2022).
High-resolution image synthesis with latent diffusion models.
\textit{IEEE/CVF Conference on Computer Vision and Pattern Recognition}, 10684--10695.

\bibitem[Saleh y Elgammal, 2015]{saleh2015wikiart}
Saleh, B., y Elgammal, A. (2015).
Large-scale classification of fine-art paintings: Learning the right metric on the right feature.
\textit{arXiv}. \url{https://arxiv.org/abs/1505.00855}

\bibitem[Shen et al., 2020]{shen2020interfacegan}
Shen, Y., Gu, J., Tang, X., y Zhou, B. (2020).
Interpreting the latent space of GANs for semantic face editing.
\textit{IEEE/CVF Conference on Computer Vision and Pattern Recognition}, 9243--9252.

\bibitem[Song et al., 2021]{song2021ddim}
Song, J., Meng, C., y Ermon, S. (2021).
Denoising diffusion implicit models.
\textit{International Conference on Learning Representations}.

\bibitem[Wang et al., 2004]{wang2004ssim}
Wang, Z., Bovik, A.~C., Sheikh, H.~R., y Simoncelli, E.~P. (2004).
Image quality assessment: From error visibility to structural similarity.
\textit{IEEE Transactions on Image Processing, 13}(4), 600--612.

\end{thebibliography}

\end{document}
